\section{Обоснование решения математической постановки задачи}

Математическая постановка задачи была сформулирована как
задача оптимизации по целочисленной переменной $k$, где $k$ интерпретируется
как длина совместимого префикса трассы
$\tau=(\sigma_1,\sigma_2,\dots,\sigma_m)$. Допустимыми значениями переменной
оптимизации являются элементы множества $K_{\mathrm{доп}}(\mathcal{M},\tau)$,
то есть такие значения длины префикса, для которых существует поведение модели,
одновременно допустимое для системы переходов $\mathcal{M}$ и согласованное с
соответствующим префиксом трассы. Требуется найти максимальное значение
показателя качества $q(k)$ на данном множестве.

Поскольку по определению
\[
q(k)=\frac{k}{m},
\qquad m \ge 1,
\]
а функция $q$ строго возрастает по $k$ на множестве $\{0,\dots,m\}$, задача
максимизации $q(k)$ эквивалентна задаче нахождения максимального элемента
множества $K_{\mathrm{доп}}(\mathcal{M},\tau)$. Следовательно, задача сводится
к определению наибольшей длины префикса трассы, совместимого с моделью.

\subsection{Основная идея решения}

Из формулировки задачи непосредственно следует естественная идея её решения.
Поскольку трасса состоит из конечного числа наблюдений, поиск максимального
совместимого префикса можно выполнять последовательно, по мере увеличения его
длины.

Для этого необходимо:
\begin{enumerate}
    \item определить множество состояний, достижимых после согласования первого
    наблюдения $\sigma_1$;
    \item затем определить множество состояний, достижимых после согласования
    первых двух наблюдений $(\sigma_1,\sigma_2)$;
    \item продолжать этот процесс вплоть до последнего наблюдения $\sigma_m$;
    \item найти наибольший индекс $k$, для которого после согласования первых
    $k$ наблюдений остаётся хотя бы одно достижимое состояние модели.
\end{enumerate}

Если после обработки всех $m$ наблюдений множество достижимых состояний
непусто, то полная трасса совместима с моделью. Если же на некотором шаге
множество достижимых состояний становится пустым, то дальнейшее продолжение
невозможно, а значит, совместим лишь некоторый собственный префикс трассы.

Тем самым задача нахождения оптимального значения показателя качества сводится
к последовательному вычислению достижимых множеств состояний, по которым затем
восстанавливается множество допустимых длин
$K_{\mathrm{доп}}(\mathcal{M},\tau)$ и его максимальный элемент.

\subsection{Рекуррентное построение множеств $D_i$}

Для формализации указанной идеи введём последовательность множеств
\[
D_0, D_1, \dots, D_m \subseteq X,
\]
где каждое множество $D_i$ интерпретируется как множество всех состояний
модели, достижимых после согласования первых $i$ наблюдений трассы.

Поскольку трасса не содержит отдельного наблюдения начального состояния,
исходное множество определяется только множеством начальных состояний модели.

\begin{definition}\label{def:d0}
Начальное множество достижимых согласованных состояний определяется равенством
\begin{equation}\label{f14}
D_0 = X_0.
\end{equation}
\end{definition}

Это определение означает, что до обработки трассы допустимыми считаются все
начальные состояния модели.

Далее, для каждого $i \in \{1,\dots,m\}$ определим множество $D_i$ рекуррентно.

\begin{definition}\label{def:di}
Для каждого $i \in \{1,\dots,m\}$ множество достижимых согласованных состояний
после обработки первых $i$ наблюдений определяется равенством
\begin{equation}\label{f15}
D_i
=
\left\{
y \in X
:
\exists x \in D_{i-1} : (x,y) \in R \land C(\sigma_i,x,y)=1
\right\}.
\end{equation}
\end{definition}

Смысл этого определения состоит в следующем. Состояние $y$ включается в
множество $D_i$ тогда и только тогда, когда существует хотя бы одно состояние
$x \in D_{i-1}$, уже достижимое после согласования первых $i-1$ наблюдений,
такое, что переход из $x$ в $y$:
\begin{itemize}
    \item является допустимым переходом модели, то есть $(x,y) \in R$;
    \item согласован с $i$-м наблюдением $\sigma_i$, то есть
    $C(\sigma_i,x,y)=1$.
\end{itemize}

Следовательно, определение \ref{def:di} отражает принцип пошагового
распространения ограничений: каждое следующее наблюдение применяется только к
тем состояниям, которые уже достижимы после обработки предыдущего префикса
трассы. Это означает, что множество $D_0$ описывает все возможные исходные
положения системы, множество $D_1$ описывает все состояния, в которые система
может перейти после одного шага, согласованного с $\sigma_1$ и т.д.

Для обоснования корректности предложенного решения необходимо доказать, что
множество $D_i$ действительно содержит в точности те состояния, которые
достижимы после согласования первых $i$ наблюдений трассы.

\begin{proposition}\label{prop:Di}
Для любого индекса $i \in \{0,\dots,m\}$ и любого состояния $x \in X$
следующие утверждения эквивалентны:
\begin{enumerate}
    \item $x \in D_i$;
    \item существует последовательность состояний
     \[
     \begin{aligned}
     &\quad(x_0,x_1,\dots,x_i)\in X^{i+1} :\\
     &\quad x_0\in X_0,\\
     &\quad x_i = x,\\
     &\quad\forall k\in\{0,\dots,i-1\}:\ (x_k,x_{k+1})\in R,\\
     &\quad\forall k\in\{1,\dots,i\}:\ C(\sigma_k,x_{k-1},x_k)=1.
     \end{aligned}
     \]
\end{enumerate}
\end{proposition}

Иначе говоря, состояние $x$ принадлежит множеству $D_i$ тогда и только тогда,
когда существует путь длины $i$, начинающийся в допустимом начальном состоянии,
допустимый для модели и согласованный с первыми $i$ наблюдениями трассы,
который оканчивается в состоянии $x$.

\begin{proof}
Доказательство проводится по индукции по $i$.

\textbf{База индукции.} Пусть $i=0$. Тогда по определению \ref{def:d0} $D_0 =
X_0$. Следовательно, для любого $x \in X$ условие $x \in D_0$ эквивалентно
условию $x \in X_0$. Это, в свою очередь, эквивалентно существованию
последовательности длины $1$, состоящей из единственного состояния $(x)$ и
удовлетворяющей условиям пункта б. База доказана.

\textbf{Переход индукции.}
Пусть утверждение верно для некоторого $i-1 \in \{0,\dots,m-1\}$. Докажем его
для индекса $i$.

\medskip

\noindent
\emph{Необходимость.}
Пусть $x \in D_i$. Тогда по определению \ref{def:di}
\[
\exists y \in D_{i-1} : (y,x) \in R \land C(\sigma_i,y,x)=1.
\]
По предположению индукции из условия $y \in D_{i-1}$ следует существование
последовательности состояний $(x_0,x_1,\dots,x_{i-1}) \in X^i$,
для которой выполнены условия
\[
\begin{aligned}
&x_0 \in X_0, \\
&x_{i-1} = y, \\
&\forall k \in \{0,\dots,i-2\} : (x_k,x_{k+1}) \in R, \\
&\forall k \in \{1,\dots,i-1\} : C(\sigma_k,x_{k-1},x_k)=1.
\end{aligned}
\]
Добавляя к этой последовательности состояние $x$, получаем последовательность
$(x_0,x_1,\dots,x_{i-1},x) \in X^{i+1}$, для которой дополнительно выполнены
условия
\[
(x_{i-1},x) \in R,
\qquad
C(\sigma_i,x_{i-1},x)=1.
\]
Следовательно, существует последовательность состояний, удовлетворяющая
условиям пункта б.

\noindent
\emph{Достаточность.}
Пусть существует последовательность $(x_0,x_1,\dots,x_i) \in X^{i+1}$, для
которой выполнены условия
\[
\begin{aligned}
&x_0 \in X_0, \\
&x_i = x, \\
&\forall k \in \{0,\dots,i-1\} : (x_k,x_{k+1}) \in R, \\
&\forall k \in \{1,\dots,i\} : C(\sigma_k,x_{k-1},x_k)=1.
\end{aligned}
\]
Тогда её префикс $(x_0,x_1,\dots,x_{i-1})$ удовлетворяет условиям пункта б для
индекса $i-1$, а потому по предположению индукции $x_{i-1} \in D_{i-1}$. Из
условий
\[
(x_{i-1},x_i) \in R
\qquad\text{и}\qquad
C(\sigma_i,x_{i-1},x_i)=1
\]
и определения \ref{def:di} получаем $x_i \in D_i$. Так как $x_i = x$, то $x \in
D_i$.

Тем самым обе импликации доказаны, следовательно, утверждения а и б
эквивалентны. По принципу математической индукции утверждение доказано для
всех $i \in \{0,\dots,m\}$.
\end{proof}

\subsection{Связь множеств $D_i$ и множества $K_{\mathrm{доп}}(\mathcal{M},\tau)$}

Поскольку итоговая математическая постановка сформулирована через множество
допустимых длин согласованных префиксов $K_{\mathrm{доп}}(\mathcal{M},\tau)$,
необходимо установить его явную связь с рекуррентно построенными множествами
$D_i$.

\begin{proposition}\label{prop:K}
Справедливо равенство
\begin{equation}\label{f16}
K_{\mathrm{доп}}(\mathcal{M},\tau)
=
\left\{
i \in \{0,\dots,m\} : D_i \neq \varnothing
\right\}.
\end{equation}
\end{proposition}

\begin{proof}
По определению \ref{def:K} число $i$ принадлежит множеству
$K_{\mathrm{доп}}(\mathcal{M},\tau)$ тогда и только тогда, когда
\[
\mathrm{B}_i(\mathcal{M}) \cap \mathrm{M}_i(\tau) \neq \varnothing.
\]
Последнее означает существование последовательности состояний
$(x_0,\dots,x_i)\in X^{i+1}$, для которой выполнены условия
\[
\begin{aligned}
&\quad x_0\in X_0,\\
&\quad x_i = x,\\
&\quad\forall k\in\{0,\dots,i-1\}:\ (x_k,x_{k+1})\in R,\\
&\quad\forall k\in\{1,\dots,i\}:\ C(\sigma_k,x_{k-1},x_k)=1.
\end{aligned}
\]

Согласно утверждению \ref{prop:Di}, это эквивалентно условию
$D_i \neq \varnothing$.
Следовательно,
\[
i \in K_{\mathrm{доп}}(\mathcal{M},\tau)
\iff
D_i \neq \varnothing,
\]
что и доказывает равенство \eqref{f16}.
\end{proof}

\subsection{Нахождение оптимального значения показателя качества}

По определению \ref{def:k_max} $k^*(\mathcal{M},\tau) = \max
K_{\mathrm{доп}}(\mathcal{M},\tau)$. С учётом утверждения \ref{prop:K} получаем
\begin{equation}\label{f17}
k^*(\mathcal{M},\tau)
=
\max
\left\{
i \in \{0,\dots,m\} : D_i \neq \varnothing
\right\}.
\end{equation}

Следовательно, решение оптимизационной задачи полностью сводится к
последовательному вычислению множеств $D_i$ и последующему поиску наибольшего
индекса, для которого соответствующее множество непусто.

\begin{theorem}\label{th:Q}
Оптимальное значение показателя качества согласованности трассы с моделью
определяется формулой
\[
Q(\mathcal{M},\tau)
=
\frac{1}{m}
\max
\left\{
i \in \{0,\dots,m\} : D_i \neq \varnothing
\right\}.
\]
\end{theorem}

\begin{proof}
По определению \ref{def:k_max}
$k^*(\mathcal{M},\tau)=\max K_{\mathrm{доп}}(\mathcal{M},\tau)$.
По утверждению \ref{prop:K}
$K_{\mathrm{доп}}(\mathcal{M},\tau)
=
\left\{
i \in \{0,\dots,m\} : D_i \neq \varnothing
\right\}.
$
Следовательно,
\[
k^*(\mathcal{M},\tau)
=
\max
\left\{
i \in \{0,\dots,m\} : D_i \neq \varnothing
\right\}.
\]
Остаётся воспользоваться определением показателя качества:
\[
Q(\mathcal{M},\tau)=\frac{k^*(\mathcal{M},\tau)}{m}.
\]
Теорема доказана.
\end{proof}

\subsection{Критерий полной совместимости трассы с моделью}

Из доказанных утверждений непосредственно вытекает критерий решения связанной
задачи разрешимости.

\begin{theorem}\label{th:decision}
Для трассы программы $\tau=(\sigma_1,\sigma_2,\dots,\sigma_m)$ следующие
утверждения эквивалентны:
\begin{enumerate}
    \item трасса $\tau$ полностью совместима с моделью $\mathcal{M}$;
    \item $m \in K_{\mathrm{доп}}(\mathcal{M},\tau)$;
    \item $D_m \neq \varnothing$;
    \item $Q(\mathcal{M},\tau)=1$.
\end{enumerate}
\end{theorem}

\begin{proof}
Эквивалентность утверждений 1 и 2 непосредственно следует из определения полной
совместимости трассы с моделью и определения \ref{def:K}.

Эквивалентность утверждений 2 и 3 следует из утверждения \ref{prop:K}, согласно
которому
\[
K_{\mathrm{доп}}(\mathcal{M},\tau)
=
\left\{
i \in \{0,\dots,m\} : D_i \neq \varnothing
\right\}.
\]

Покажем эквивалентность утверждений 2 и 4. Если $m \in
K_{\mathrm{доп}}(\mathcal{M},\tau)$, то по определению \ref{def:k_max}
$k^*(\mathcal{M},\tau)=m$, так как всегда $k^*(\mathcal{M},\tau)\le m$.
Следовательно,
\[
Q(\mathcal{M},\tau)=\frac{k^*(\mathcal{M},\tau)}{m}=1.
\]

Обратно, если
$Q(\mathcal{M},\tau)=1$,
то по определению показателя качества
\[
\frac{k^*(\mathcal{M},\tau)}{m}=1,
\]
откуда $k^*(\mathcal{M},\tau)=m$. Следовательно, $m \in
K_{\mathrm{доп}}(\mathcal{M},\tau)$.

Тем самым все четыре утверждения эквивалентны.
\end{proof}

Таким образом, задача разрешимости оказывается частным случаем оптимизационной
постановки: необходимо определить, достигается ли наилучшее значение
показателя качества.

\subsection{Алгоритм решения}

Из построения множеств $D_i$ непосредственно следует алгоритм решения как
оптимизационной задачи, так и связанной с ней задачи разрешимости. Алгоритм
последовательно вычисляет множества достижимых согласованных состояний,
тем самым неявно определяя множество допустимых длин
$K_{\mathrm{доп}}(\mathcal{M},\tau)$, после чего находит его максимальный
элемент $k^*$ и вычисляет значение показателя качества.

\begin{algorithm}[htp!]
    \SetAlgoLined
    \KwData{система переходов $\mathcal{M} = (X, X_0, R)$; трасса $\tau = (\sigma_1,\dots,\sigma_m) \in \Sigma^m$; предикат согласованности $C : \Sigma \times X \times X \to \{0,1\}$}
    \KwResult{значение $k^*(\mathcal{M},\tau)$, значение $Q(\mathcal{M},\tau)$ и ответ о полной совместимости трассы}

    $D_0 \gets X_0$\;
    $k^* \gets 0$\;

    \ForEach{$i \in \{1,\dots,m\}$}{
        $D_i \gets \left\{
            y \in X
            :
            \exists x \in D_{i-1} : (x,y) \in R \land C(\sigma_i,x,y)=1
        \right\}$\;

        \If{$D_i = \varnothing$}{
            \textbf{прервать цикл}\;
        }

        $k^* \gets i$\;
    }

    $Q(\mathcal{M},\tau) \gets \frac{k^*}{m}$\;

    \If{$k^* = m$}{
        \Return{$k^*$, $Q(\mathcal{M},\tau)$, $\tau$ полностью совместима с $\mathcal{M}$}\;
    }
    \Return{$k^*$, $Q(\mathcal{M},\tau)$, $\tau$ не полностью совместима с $\mathcal{M}$}\;
    \caption{Нахождение максимального согласованного префикса трассы}
    \label{alg:trace-quality}
\end{algorithm}

Данный алгоритм является корректным в силу теорем \ref{th:Q} и
\ref{th:decision}. Если множество состояний $X$ конечно, то алгоритм является
конструктивно осуществимым, так как каждое множество $D_i$ вычисляется за
конечное число шагов.

Особый интерес представляет случай, когда $Q(\mathcal{M},\tau) < 1$. Тогда
полная трасса не совместима с моделью. Однако это не означает, что не
существует ни одного совместимого префикса трассы. Напротив, существует
максимальный индекс
\[
k^*(\mathcal{M},\tau)
=
\max
\left\{
i \in \{0,\dots,m\} : D_i \neq \varnothing
\right\},
\]
после которого дальнейшее продолжение становится невозможным.

Смысл числа $k^*(\mathcal{M},\tau)$ состоит в том, что первые
$k^*(\mathcal{M},\tau)$ наблюдений трассы ещё могут быть объяснены некоторым
допустимым поведением модели, тогда как наблюдение
$\sigma_{k^*(\mathcal{M},\tau)+1}$ уже не может быть согласовано ни с одним
продолжением. Тем самым значение $k^*(\mathcal{M},\tau)$ указывает на первый
шаг трассы, на котором возникает расхождение между реализацией и моделью, а
значение $Q(\mathcal{M},\tau)$ характеризует степень близости трассы к полному
согласованию.

Это наблюдение важно не только теоретически, но и практически, поскольку
позволяет использовать предложенный метод не только для принятия или
отклонения трассы в целом, но и для количественной оценки степени
несоответствия и для локализации шага, на котором оно возникает.
